\documentclass{article}
\usepackage{amsmath,amsthm,amssymb}
\usepackage[margin=1in]{geometry}
\usepackage{url}

\DeclareMathOperator{\area}{area}
\DeclareMathOperator{\dinv}{dinv}
\DeclareMathOperator{\RR}{RR}

\newtheorem{theorem}{Theorem}
\newtheorem{conjecture}{Conjecture}
\newtheorem{remark}{Remark}

\title{Dyck Skeleton--Tableau Formulas}
\author{}
\date{}

\begin{document}
\maketitle

\section{Overview}

This note records a two-column skeleton--tableau formula for the
\(q,t\)-Catalan polynomial.  In the classical Dyck case the formula is a
proved theorem of the 2026 paper \emph{Dyck Symmetric Functions and
Applications to \(q,t\)-Catalan Polynomials}.  The same expression has a
natural rational analogue in the family
\[
  r=s\tau+1,
\]
where \(s\) is the path length and \(\tau\) is a nonnegative integer step
parameter.  For \(\tau>1\) this rational analogue is conjectural, but it has
been checked in finite ranges by comparing the direct \(q,t\)-coefficient
polynomial with the skeleton--tableau expansion.

The formula is useful because it separates a Dyck path into two pieces.  The
first piece is a skeleton \(F\), which records a rigid part of the path.  The
second piece is an at-most-two-column Dyck tableau \(P\), whose shape produces
a two-variable Schur factor.  Thus the formula rewrites a sum over individual
Dyck paths as a sum over pairs \((F,P)\).

Throughout the note, \(q\) records area and \(t\) records dinv.  The step
parameter is written \(\tau\) to avoid confusing it with the variable \(t\).

\section{The classical formula}

A classical Dyck sequence of length \(n\) is a sequence
\[
  D=(D_0,\ldots,D_{n-1})
\]
of nonnegative integers such that \(D_0=0\) and
\[
  D_{i+1}\le D_i+1
  \qquad(0\le i<n-1).
\]
Its area is
\[
  \area(D)=\sum_i D_i,
\]
and its classical dinv statistic is
\[
  \dinv(D)=\#\{(i,j):0\le i<j<n,
      \ D_i=D_j\text{ or }D_i=D_j+1\}.
\]
The ordinary \(q,t\)-Catalan polynomial in this convention is
\[
  C_n(q,t)=\sum_D q^{\area(D)}t^{\dinv(D)},
\]
where \(D\) ranges over all Dyck sequences of length \(n\).

The formula below uses Dyck \(m\)-skeletons and Dyck tableaux.  Informally, a
Dyck \(m\)-skeleton is a Dyck sequence whose final entry is its maximum
\(m\), and which has no nonfinal entry that can be extracted by the skeleton
reduction rule.  A Dyck tableau is a left-aligned tableau whose rows are dual
Dyck sequences and whose columns are affine Dyck sequences.  The formula only
uses tableaux with at most two columns and with entries in \([0,m-1]\).  If
\(P\) is such a tableau, \(\lambda(P)\) denotes its partition shape and
\(\RR(P)\) denotes its row-reading word, read from bottom row to top row and
left to right inside each row.

\begin{theorem}[Hawkes, two-column tableau formula]
For every \(n\ge1\),
\[
  C_n(q,t)
  =
  \sum_{(F,P)}
  q^{\area(F:\RR(P))}
  t^{\dinv(F:\RR(P))-|P|}
  s_{\lambda(P)'}(q,t).
\]
The sum is over all pairs \((F,P)\) such that \(F\) is a Dyck
\(m\)-skeleton for some \(m\ge0\), \(P\) is an at-most-two-column Dyck
tableau with entries in \([0,m-1]\), and
\[
  |F|+|\RR(P)|=n.
\]
Here \(\lambda(P)'\) is the conjugate partition and
\(s_{\lambda(P)'}(q,t)\) is the Schur function in the two variables \(q,t\).
\end{theorem}

The Schur factor comes from the proof.  The 2026 paper constructs a chain of
area- and dinv-preserving correspondences from Dyck sequences to triples
\((F,P,Q)\), where \(F\) is the skeleton, \(P\) is the two-column Dyck
tableau, and \(Q\) is a binary reverse semistandard tableau of the same shape
as \(P\).  After \((F,P)\) is fixed, summing over all possible binary
recording tableaux \(Q\) gives the two-variable Schur function on the
conjugate shape.  The factor \(t^{-|P|}\) appears because a cell of \(Q\)
contributes either a \(q\)-weight or a compensating negative \(t\)-shift before
this sum is rewritten as a Schur polynomial.

\section{The rational \(r=s\tau+1\) model}

Now fix integers \(s\ge1\) and \(\tau\ge0\).  A normalized \(\tau\)-Dyck
sequence of length \(s\) is a sequence
\[
  D=(D_0,\ldots,D_{s-1})
\]
of nonnegative integers such that
\[
  D_0=0,
  \qquad
  D_{i+1}\le D_i+\tau
  \qquad(0\le i<s-1).
\]
For any finite integer word \(x=(x_0,\ldots,x_{\ell-1})\), define
\[
  \dinv_\tau(x)
  =
  \sum_{0\le i<j<\ell} d_\tau(x_i,x_j),
\]
where
\[
  d_\tau(a,b)=
  \begin{cases}
    \max(0,a+\tau-b),& a\le b,\\
    \max(0,b+1+\tau-a),& a>b.
  \end{cases}
\]
The rational \(q,t\)-Catalan polynomial in this model is
\[
  C_{s,\tau}(q,t)
  =
  \sum_D q^{\area(D)}t^{\dinv_\tau(D)},
\]
where \(D\) ranges over normalized \(\tau\)-Dyck sequences of length \(s\).
For \(\tau=1\) this is the classical Dyck case with \(n=s\).

The rational skeleton and tableau definitions parallel the classical ones.  An
entry \(D_j=e\) of a normalized \(\tau\)-Dyck sequence is called extractable if
\(e>0\), exactly one earlier entry lies in the predecessor interval
\[
  [\max(0,e-\tau),e),
\]
and deleting \(D_j\) preserves the adjacent \(\tau\)-Dyck inequality.  If
\(0<j<s-1\), the deletion condition is
\[
  D_{j+1}\le D_{j-1}+\tau.
\]
If \(j=s-1\), there is no new adjacent pair to check.  A rational
\(m\)-skeleton is a normalized \(\tau\)-Dyck sequence whose final entry is its
maximum \(m\), and which has no nonfinal extractable entry.  A final
extractable entry is allowed.

A rational Dyck tableau of step \(\tau\) is a left-aligned tableau whose row
lengths form a partition, whose rows are dual \(\tau\)-Dyck sequences
\[
  P_i[j+1]>P_i[j]+\tau,
\]
and whose columns, read from bottom to top, are affine \(\tau\)-Dyck
sequences.  Equivalently, if row \(i\) is immediately above row \(i+1\), then
\[
  P_i[j]\le P_{i+1}[j]+\tau
\]
whenever both entries exist.  Its row-reading word \(\RR(P)\) is read from
bottom row to top row, left to right within each row.

\begin{conjecture}[Rational two-column skeleton--tableau formula]
For \(\tau\ge0\) and \(s\ge1\),
\[
  C_{s,\tau}(q,t)
  =
  \sum_{(F,P)}
  q^{\area(F:\RR(P))}
  t^{\dinv_\tau(F:\RR(P))-|P|}
  s_{\lambda(P)'}(q,t).
\]
The sum is over all pairs \((F,P)\) such that \(F\) is a rational
\(m\)-skeleton of step \(\tau\) for some \(m\ge0\), \(P\) is an
at-most-two-column rational Dyck tableau of step \(\tau\) with entries in
\([0,m-1]\), and
\[
  |F|+|\RR(P)|=s.
\]
\end{conjecture}

For \(\tau=1\), this is exactly the classical theorem above.  For \(\tau=0\),
the model is degenerate: the only normalized \(\tau\)-Dyck sequence of each
length is the all-zero sequence, and the formula reduces to the single empty
tableau contribution.  For \(\tau>1\), the displayed identity is a conjecture.
The code accompanying this item checks it by direct finite enumeration.

\section{Computational checks}

The file \texttt{code.py} is a readable reference checker for the rational
formula.  For each requested pair \((\tau,s)\), it builds two coefficient
Dictionaries indexed by \((\area,\dinv)\):
\begin{itemize}
  \item the direct dictionary, obtained by enumerating all normalized
        \(\tau\)-Dyck sequences of length \(s\);
  \item the formula dictionary, obtained by enumerating rational
        \(m\)-skeletons, at-most-two-column rational Dyck tableaux, and the
        two-variable Schur expansion of \(s_{\lambda(P)'}(q,t)\).
\end{itemize}
The check passes exactly when these two dictionaries agree for every exponent
pair.

The default command is deliberately small:
\[
  \texttt{python code.py}
\]
It checks \(\tau=2\) and \(1\le s\le5\).  Larger cases were verified for this
item with optimized code.  The recorded verified ranges are
\[
\begin{array}{c|c}
\tau & \text{verified lengths}\\
\hline
2 & 1\le s\le14\\
3 & 1\le s\le12\\
4 & 1\le s\le10.
\end{array}
\]
Those optimized runs compare the same two coefficient dictionaries, but use a
more efficient implementation than the readable file included here.

The computational evidence does not prove the conjecture for all \(s\) when
\(\tau>1\).  It does, however, give a precise finite formulation: for each
checked \((\tau,s)\), every coefficient of the direct rational
\(q,t\)-Catalan polynomial agrees with the corresponding coefficient predicted
by the rational skeleton--tableau expansion.

\section{Status}

The classical two-column skeleton--tableau formula is theorem-level material:
it is Theorem 4.20 of the 2026 paper.  The rational \(r=s\tau+1\) formula is
known in the classical specialization \(\tau=1\) and is tautological in the
degenerate specialization \(\tau=0\).  For \(\tau>1\), it should be read as a
conjectural extension supported by the finite checks listed above.  A general
proof, or a conceptual bijection explaining why the same two-column
skeleton--tableau structure survives for \(\tau>1\), remains open.

\begin{thebibliography}{9}

\bibitem{Hawkes2026DyckSymmetric}
Graham Hawkes,
\emph{Dyck Symmetric Functions and Applications to \(q,t\)-Catalan Polynomials},
arXiv:2605.13003, 2026.
\[
  \texttt{https://arxiv.org/abs/2605.13003}
\]

\end{thebibliography}

\end{document}
